\chapter{Identification, counterfactuals and effect sets}
\label{exec17:s3}

A bank can know that an offer increases average completion and still be unable
to say which completed applications the offer caused. It can identify the
effect of sending an encouragement without identifying the effect of making
every customer receive the service. It can also possess a valid identifying
formula whose numerical evaluation is too fragile to support the proposed
decision. These are different limits on what a causal product can say. The
choice between software components depends on preserving those limits through
the calculation.

This chapter develops that choice around a binary servicing action and a
binary successful outcome. The examples are project-authored probability
laws, evaluated exactly, with a stipulated value of USD 100 per success and
USD 12 per action received. They are neither estimates of bank behavior nor
measurements of upstream software. Their role is to make a proposal testable:
an implementation can be asked to preserve the target, recover the
appropriate answer or set, and explain a refusal without confusing a
scientific obstacle with a computational one.~\citep{exec17_s3_exact}

\section{A causal question is relative to information and a model class}
\label{exec17:s3-target}

Let $A\in\{0,1\}$ denote receipt of a specified intervention, $Y\in\{0,1\}$
the outcome at a fixed horizon, and $C$ information measured before receipt.
The intervention includes its delivery version: sending a message and
completing a counseling session are different actions even if both are
recorded as contact. Write $Y(a)$ for the response under receipt $a$,
$P$ for the observable population law and $\mathcal M$ for the admitted
causal models. A common target is
\[
 \theta(M)=E_M\{Y(1)-Y(0)\},\qquad M\in\mathcal M.
\]
Consistency links the observed outcome to the potential response under
the action actually received. Assumptions about delivery versions,
interference, information availability and the horizon determine whether
that link describes the banking question. A graph supplies restrictions
on mechanisms and omitted causes; it does not certify these operational
premises.~\citep{pearl2009,swig2013}

For a fixed observable law, define the compatible model class and effect set by
\[
 \mathcal M(P)=\{M\in\mathcal M:P_M=P\},\qquad
 \Theta(P)=\{\theta(M):M\in\mathcal M(P)\}.
\]
The target is point identified at $P$ if this set is a singleton. A graphical
identification theorem usually establishes that a functional $T$ works for
every admitted compatible law, so that $\theta(M)=T(P_M)$ throughout its
stated domain. Special numerical coincidences at one law are weaker evidence.
Conversely, finding two admissible models with the same $P$ and different
$\theta$ disproves point identification for that law and class. Estimating
$P$ more accurately does not distinguish such models. A smaller scientifically
defensible class or additional intervention information may
do so.~\citep{id2006,pearl2009}

The distinction matters when evaluating a proposal. A held-out likelihood
score assesses a distributional prediction. It cannot establish uniqueness
of a causal answer across models that share that distribution. An optimizer
that returns one fitted model has selected one element; it has not shown that
all compatible elements agree. A prior can produce a posterior over models,
but its concentration has to be interpreted with the restrictions and prior
that supplied the missing information.

There is another boundary: a formula may involve $P(Y\mid A=a,C=c)$ where
$P(A=a\mid C=c)=0$. The observed law then does not supply that conditional
response. A smooth extrapolating regression can return a value, but the value
uses additional restrictions. The proposal should declare whether it restricts
the target population, changes the action, acquires information or adopts
an explicit extrapolation model. A numerical nonzero fitted propensity
does not create observed support.

\section{Adjustment changes the comparison, not the action definition}
\label{exec17:s3-adjust}

Suppose consistency, conditional exchangeability
$Y(a)\mathrel{\perp}A\mid C$ and positive receipt probabilities hold in the
target population. Iterated expectation then gives
\begin{align}
 E\{Y(a)\}
 &=\sum_c E\{Y(a)\mid C=c\}P(C=c) \notag\\
 &=\sum_c E\{Y\mid A=a,C=c\}P(C=c).
 \label{exec17:s3-adjustment}
\end{align}
The first equality standardizes to the chosen population; the second uses
exchangeability and consistency on the supported cells. If the sponsor wants
the treated population, replacing $P(C=c)$ by $P(C=c\mid A=1)$ changes the
target. It is not an innocuous estimation option.

An exact generated law makes the cost of ignoring that distinction visible.
Let $P(C=0)=P(C=1)=1/2$ and stipulate
\[
\begin{array}{c|cc|cc}
 &P(A=1\mid C)&P(A=0\mid C)
 &P(Y=1\mid A=0,C)&P(Y=1\mid A=1,C)\\ \hline
 C=0&1/5&4/5&1/10&1/5\\
 C=1&4/5&1/5&3/5&7/10
\end{array}
\]
This is a stipulated exchangeable-within-$C$ construction, rather than
evidence that a real baseline variable removes every common cause.
Both strata have a receipt effect of $1/10$. Standardization yields
$E\{Y(0)\}=7/20$ and $E\{Y(1)\}=9/20$.
The observed receipt groups have different compositions:
\[
 P(C=1\mid A=1)=4/5,\qquad P(C=1\mid A=0)=1/5.
\]
Their outcome means are $3/5$ and $1/5$, producing a raw contrast of
$2/5$. All cells are positive, so the problem is not an absent observation
cell. Receipt selection overweights the high-success stratum in the
treated group.~\citep{exec17_s3_exact}

For universal receipt versus no receipt, define incremental expected
net value $V=100\theta-12$ in USD per eligible customer. The causal contrast
implies $V=-2$; substituting the association implies $V=28$.
A product that chooses an accurate predictive model but compares its
predictions over the wrong receipt groups can therefore recommend the
opposite action. The USD 30 discrepancy here is exact under the generated
law, not a forecast of attainable financial improvement.

The relevant diagnostics follow the derivation. Check treatment versions
and baseline timing; inspect supported cells and overlap; state the
standardization weights; and compare the proposed graph against substantive
common causes and selection. Balance after weighting can diagnose a
particular imbalance in measured variables. It does not demonstrate
exchangeability with respect to unmeasured variables.

Adjustment can also answer the wrong causal question when it conditions
on a mediator. Consider the separate linear structural example
\[
 M=A+\varepsilon_M,\qquad
 Y=2A+3M+\varepsilon_Y,
\]
with randomized $A$ and independent zero-mean errors. Substitution gives
$E\{Y\mid do(A=a)\}=5a$ up to the zero baseline intercept. The total
mean effect of changing $A$ from zero to one is $5$. Holding $M$ fixed
by intervention instead gives a controlled direct effect of $2$.
These are continuous illustrative responses, not the binary success law
above. Including the mediator has changed the intervention comparison.
In a bank, completion assistance and the downstream information it
elicits may both contribute to the total effect; controlling away the
information channel is appropriate only if that direct-effect question
is actually intended.~\citep{exec17_s3_exact,pearl2009}

\section{Completeness must retain the procedure's domain}

Perkovi\'c and colleagues characterize adjustment in DAGs, CPDAGs,
MAGs and PAGs using amenability, exclusion of a forbidden set and blocking
of proper definite-status noncausal paths. Their Theorem 5 gives an
if-and-only-if result for adjustment; Theorem 7 supplies the equivalent
proper-back-door separation condition. The MAG and PAG treatment in this
source excludes selection bias. These conditions require the graph's
endpoint semantics, not merely its adjacency
matrix.~\citep{perkovic2018}

Completeness here means that the criterion characterizes the existence
of adjustment sets within the stated classes and definition. It does
not mean that failure to find such a set proves failure of every
identification method. The same source explicitly distinguishes that
claim from identification by other functionals. An adapter should
therefore report the attempted route and its domain with its result.
No valid adjustment set in the admitted graph leaves a different
question open from nonidentification in the admitted model class.
A search that stops after small candidate sets is weaker again,
because it may not have exhausted the adjustment
problem.~\citep{perkovic2018}

For a supplied CPDAG or PAG, a successful adjustment set must be valid
under the relevant equivalence-class interpretation. Selecting one
convenient orientation can create an answer that does not survive the
remaining graphs. It is reasonable to ask a domain expert for the
orientation of a particular edge, but that additional restriction must
remain visible in the scientific model. The resulting effect is
conditional on it.

This motivates three substantive alternatives when adjustment fails.
One can seek a more general identifying functional under the same
information, retain a set of effects compatible with the information,
or change the information through an admissible experiment or
measurement. Choosing a convenient regression is a fourth computational
action, but it only answers the causal question if its added
restrictions are justified. The banking consequence should be assessed
before choosing among these routes: a set lying entirely below the
receipt cost threshold can already rule out an action.

\section{Joint and conditional intervention questions need different routes}
\label{exec17:s3-idc}

For a semi-Markovian causal graph, ID identifies unconditional
intervention distributions or returns a graphical nonidentification
witness in its formal domain. The district decomposition and hedge
criterion are scientific statements about that model class.
They are not properties of every exception a software implementation
can throw. Shpitser and Pearl's conditional algorithm, IDC, first
uses action--observation exchange where the graphical condition permits
it. After the maximal eligible exchange, it invokes the appropriate
joint identification problem and normalizes. Their Theorems 6--9
establish the conditional reduction, soundness, completeness and
do-calculus result for the stated
queries.~\citep{id2006}

The order matters. In general it is valid to normalize an identified
$P(Y,Z\mid do(A=a))$ to obtain $P(Y\mid do(A=a),Z=z)$
where the denominator is positive. It is invalid to infer that
the conditional target is nonidentified merely because this particular
joint target is nonidentified. The conditional may cancel a component
that the joint requires. A component that supports only joint-first
normalization is consequently a conservative partial route to
conditional queries.

Consider the graph with $A\to Z\to Y$ and $A\leftrightarrow Z$,
with no hidden common cause involving $Y$. Let $U$ be a fair bit,
$E$ an independent Bernoulli$(1/4)$ bit, and let the outcome be generated
from an independent random input with
\[
 P(Y=1\mid Z=z)=1/5+(3/5)z.
\]
Two structural models compatible with the graph are
\[
\begin{array}{c|cc}
 & A&Z\\ \hline
 M_0&U&U\mathbin{\mathrm{XOR}}E\\
 M_1&U&A\mathbin{\mathrm{XOR}}E
\end{array}
\]
The graph permits an arrow with zero effect in a particular model;
no faithfulness assumption is imposed on this witness class.
Observationally $A=U$, so both models give the same strictly positive
law for every triple $(A,Z,Y)$. Under $do(A=a)$, however, $U$ remains
random in $M_0$, whereas $Z$ responds to the set action in
$M_1$.~\citep{exec17_s3_exact}

The interventional means are
\[
\begin{array}{c|cc|c}
 &E\{Y\mid do(A=0)\}&E\{Y\mid do(A=1)\}&\text{contrast}\\ \hline
 M_0&1/2&1/2&0\\
 M_1&7/20&13/20&3/10 .
\end{array}
\]
Thus the observational law cannot identify the marginal effect in
this class; it also cannot identify the joint law of $(Y,Z)$ under
the action. Yet throughout the stated graph class the outcome mechanism
has only the observed parent $Z$ and independent exogenous input.
Consequently, on supported conditioning events,
\[
 P(Y=1\mid do(A=a),Z=z)=P(Y=1\mid Z=z)
                     =1/5+(3/5)z.
\]
The equality follows from the graph's mechanism restriction; agreement
of the two numerical witnesses alone would not prove identification.
All relevant conditioning events are positive in the displayed
witnesses.~\citep{id2006,exec17_s3_exact}

In a servicing interpretation, $Z$ could be a completed information
step and $Y$ a later successful resolution. Knowing the outcome response
among a specified post-action $Z$ group does not tell the bank how much
an action changes the frequency of that group. The conditional success
probability therefore does not substitute for the total-value question.
Conditioning on the observed $Z$ also does not identify the effect
within the principal stratum that would attain that $Z$ under either
action. That stratum is a different, cross-world target.

\section{A population response is not an individual counterfactual}
\label{exec17:s3-coupling}

Single-world intervention graphs make counterfactual variables under
one intervention explicit and connect their factorization to the
extended g-formula. Richardson and Robins distinguish the FFRCISTG
model from the stronger NPSEM with independent errors; single-world
and cross-world restrictions should not be silently interchanged.
Pearl's structural definition computes an individual's counterfactual
by changing the intervention while retaining the same exogenous state.
That gives a counterfactual meaning within a fully specified structural
model. It does not show that observational or experimental marginals
identify that structural specification.~\citep{swig2013,pearl2009}

The distinction is exact even for a randomized action. Suppose the
experiment identifies
\[
 P\{Y(0)=1\}=2/5,\qquad P\{Y(1)=1\}=3/5.
\]
Let $t=P\{Y(0)=1,Y(1)=1\}$. Every response-type law below has these
same two marginals:
\[
\begin{array}{c|cccc}
(Y(0),Y(1))&(0,0)&(0,1)&(1,0)&(1,1)\\ \hline
P&t&3/5-t&2/5-t&t .
\end{array}
\]
Nonnegative cell masses require and suffice for $0\le t\le2/5$.
Randomizing $A$ independently of the response type then produces
the same complete observable $(A,Y)$ law for every $t$.
Intermediate $t$ values are feasible as well as the endpoints,
so these are sharp ranges within the stipulated response-type
class.~\citep{exec17_s3_exact}

The average effect is always $1/5$, but the fraction helped is
$3/5-t\in[1/5,3/5]$ and the fraction harmed is
$2/5-t\in[0,2/5]$. At $t=0$, sixty percent benefit and forty
percent are harmed. At $t=2/5$, twenty percent benefit and nobody
is harmed. Both populations have the same average effect and
randomized outcome frequencies. Larger replications of the same
parallel-arm experiment do not select between these couplings.

The retrospective attribution question is different again. Among
treated successes, consistency and randomization imply
\begin{align*}
 P\{Y(0)=0\mid A=1,Y=1\}
 &=\frac{P\{Y(0)=0,Y(1)=1\}}{P\{Y(1)=1\}}\\
 &=\frac{3/5-t}{3/5}\in[1/3,1].
\end{align*}
One can therefore identify a positive average effect and still be
unable to assign a unique probability that a treated success was
caused by the treatment. An explanation screen that labels a
particular completion as caused by an offer would need additional
information and assumptions.

Assuming monotonic response, $Y(1)\ge Y(0)$ for every customer,
eliminates the harmful response type, forcing $t=2/5$ here.
It implies a benefit fraction $1/5$ and attribution probability
$1/3$. This is a meaningful substantive restriction: assistance
must never prevent a success that would otherwise occur. Delays,
confusion or competing actions can make it implausible in a bank.
The positive average experiment result neither proves nor requires
this restriction.

Expected net value under universal receipt is nevertheless
$100(1/5)-12=8$ USD in every coupling, under the stipulated linear
reward and uniform receipt cost. A sponsor concerned only with
that average reward can compare receipt with no receipt without
resolving $t$. A sponsor imposing a bound on the fraction harmed
cannot. The unresolved individual quantity becomes economically
material through the actual loss function or constraint, rather
than by declaring every nonidentified quantity equally blocking.

\section{Encouragement, receipt and the population are separate targets}
\label{exec17:s3-iv}

Suppose $Z$ is a randomized encouragement, $D$ receipt of assistance
and $Y$ the outcome. Let $D(z)$ be receipt under encouragement
$z$. With stable treatment definitions, assignment independence,
exclusion of a direct encouragement effect, relevance and
monotonicity $D(1)\ge D(0)$, the Wald ratio identifies the mean
receipt effect among compliers. This is the LATE interpretation
developed by Angrist, Imbens and Rubin, rather than a general
population ATE statement.~\citep{angrist1996}

The distinction follows algebraically. Exclusion lets us write
$Y(z,d)=Y(d)$. For binary receipt,
\[
 Y\{D(1)\}-Y\{D(0)\}
   =\{Y(1)-Y(0)\}\{D(1)-D(0)\}.
\]
Monotonicity makes the second factor one for compliers and zero
for never-takers and always-takers. Averaging and using
random assignment gives
\[
 \frac{E(Y\mid Z=1)-E(Y\mid Z=0)}
      {E(D\mid Z=1)-E(D\mid Z=0)}
 =E\{Y(1)-Y(0)\mid D(1)>D(0)\}.
\]
Both the nonzero denominator and the change from receipt to
encouragement are part of the argument. A large fitted first
stage does not verify exclusion or assignment
independence.~\citep{angrist1996}

Consider the authored law with principal strata and observable
response means
\[
\begin{array}{c|c|cc|cc}
 \text{stratum}&P&D(0)&D(1)&E\{Y(0)\}&E\{Y(1)\}\\ \hline
 \text{never}&3/10&0&0&1/5&\text{unknown}\\
 \text{complier}&2/5&0&1&3/10&1/2\\
 \text{always}&3/10&1&1&\text{unknown}&2/5 .
\end{array}
\]
Encouragement is an independent fair bit. Within each stratum
we can construct binary potential responses with the stated
means; product couplings provide one explicit construction.
The receipt means under $Z=0,1$ are $3/10,7/10$;
the outcome means are $3/10,19/50$. The first stage is $2/5$,
the outcome intention-to-treat effect is $2/25$, and LATE is
$1/5$.~\citep{exec17_s3_exact}

What remains unknown is not remedied by calling LATE the effect
for an average customer. Let $n=E\{Y(1)\mid\text{never}\}$
and $b=E\{Y(0)\mid\text{always}\}$. Each may range from zero
to one without changing any observable $(Z,D,Y)$ probability.
The population receipt effect is
\begin{align*}
 \theta_{\mathrm{pop}}
 &=\tfrac3{10}(n-\tfrac15)
   +\tfrac25(\tfrac12-\tfrac3{10})
   +\tfrac3{10}(\tfrac25-b)\\
 &=\tfrac7{50}+\tfrac3{10}(n-b)
   \in[-4/25,11/25].
\end{align*}
Both endpoints and every intermediate value are attainable.
The calculation gives sharp bounds for this particular observed
law under the declared binary IV class, not a claim that this
simple expression is the general IV bounds algorithm.

Universal receipt versus no receipt has net value
in $[-28,32]$ USD per person. Receipt for a complier has mean
incremental net value $8$ USD, but a bank cannot generally
identify each person's principal stratum from the observed
assignment and receipt. It cannot implement treating only
compliers merely because their average effect is
identified.~\citep{angrist1996,exec17_s3_exact}

Changing the feasible action to universal encouragement gives
a different, implementable comparison. Only the $2/5$ compliers
change receipt; with the stipulated success value and receipt
cost, incremental net value is
\[
 100(2/25)-12(2/5)=16/5\ \text{USD}.
\]
A direct encouragement cost of USD 2 reduces this to $6/5$.
Thus a proposal can support the encouragement action under its
assumptions while refusing to claim that universal receipt is
beneficial. The difference is substantive, not merely a choice
between two estimation libraries.

\section{Instrument diagnostics do not supply the missing assumptions}

The IV analysis suggests concrete review questions. Was the
encouragement randomized at a customer or cluster level? Did a
message alter outcomes through attention even when no assistance
was received? Could some customers avoid assistance because of
the encouragement, violating monotonicity? Are nonreceipt and
receipt defined consistently across encouragement groups? These
questions determine the target model before selecting a Wald,
two-stage or structural estimator.

Observed constraints may reject parts of an IV model; a fit that
satisfies them does not certify every unobserved restriction.
Likewise a numerical first-stage diagnostic assesses a feature
of the observable encouragement law. It cannot establish the
population effect among never-takers under receipt, or among
always-takers under nonreceipt. The construction above holds
the entire observed law fixed while changing those two
responses, so no diagnostic based solely on that law can
identify them.~\citep{angrist1996,exec17_s3_exact}

Numerical and inferential weakness also need distinct treatment.
With a fixed positive first-stage denominator $d$, perturbing
the outcome intention-to-treat numerator by $\delta$ changes
the ratio by $\delta/d$. A numerator perturbation of $1/1000$
changes the ratio by $1/400$ at $d=2/5$ and by $1/20$ at
$d=1/50$. At USD 100 per success these are USD $1/4$ and
USD 5 per receipt. This exact algebra explains why small
denominators demand precision. It is not a weak-instrument
confidence procedure: a random denominator and its dependence
on the numerator require an inferential argument of their
own.~\citep{exec17_s3_exact}

A conditional population extrapolation is possible only with
added information or restrictions. A new experiment that can
change receipt in previously unchanged groups supplies a
different information route. A structural effect-homogeneity
assumption supplies a different model route. Repeatedly
estimating the same LATE with greater precision supplies
neither. The proposal should compare feasibility, scientific
credibility and cost of these routes before promising that a
larger historical sample will settle the population question.

\section{Effect sets can already guide a decision}
\label{exec17:s3-bounds}

Return to the first generated law, but now suppose the analyst
receives only $(A,Y)$ and assumes consistency, with no
exchangeability or instrumental restriction. Its observed
cell probabilities are
\[
\begin{array}{c|cc}
 &Y=0&Y=1\\ \hline
 A=0&2/5&1/10\\
 A=1&1/5&3/10 .
\end{array}
\]
The treated successes reveal a mass $3/10$ with $Y(1)=1$.
Among untreated customers, a total mass $1/2$, the treated
potential response is missing. Hence
$E\{Y(1)\}\in[3/10,4/5]$. Similarly,
$E\{Y(0)\}\in[1/10,3/5]$.
The target ATE therefore lies in
\[
 [-3/10,7/10].
\]
For the lower endpoint, give every untreated customer a
missing treated response of zero and every treated customer
a missing untreated response of one. Reverse those choices
for the upper endpoint. Both completions satisfy consistency
and preserve every observed cell; mixtures attain the
intermediate values. Exact enumeration of the sixteen
cellwise extreme completions verifies the construction.
The limits are sharp for this law and model
class.~\citep{exec17_s3_exact}

The point answer $1/10$ obtained with the baseline and
exchangeability assumptions lies inside this set, but the
set does not select it. The two analyses answer questions
under different information and restrictions. More rows
from the unadjusted observational table reduce sampling
error around that table; they do not impose exchangeability.
Recovering a useful baseline is potentially valuable only
if the scientific assignment argument justifies how it
changes the compatible class.

With linear reward, mapping the effect set through
$V=100\theta-12$ gives $[-42,58]$ USD.
For a risk-neutral expected-value decision, a point
probability distribution over compatible models would
produce a mean value, but no such distribution is supplied
by the bounds. Selecting their midpoint adds an arbitrary
decision rule. Under a maximin criterion, no receipt has
value zero and universal receipt has worst value $-42$;
no receipt is preferred in this two-action problem.

Under minimax regret the calculation differs. Relative
to the best of receipt and no receipt in each compatible
world, receipt's worst regret is $42$, whereas no receipt's
worst regret is $58$. If only these deterministic actions
are admissible, minimax regret prefers receipt. Allowing a
randomized receipt fraction $r$, with no interference and
proportional costs, equalizes the endpoint regrets at
\[
 42r=58(1-r),\qquad r=29/50.
\]
The common regret is $609/25$ USD. This is a derived
decision illustration, not an ethical recommendation to
randomize a banking service. Operational, equity and
capacity constraints may exclude that action. It shows
why a causal product must obtain the sponsor's decision
criterion rather than turn an ambiguous effect into an
unexplained binary recommendation.~\citep{exec17_s3_decision}

\section{Identification sets, confidence sets and numerical envelopes}
\label{exec17:s3-numerical}

The preceding ranges describe variation among causal models
with the same exact observable law. They are not confidence
intervals from sampled data. If a separate justified
procedure constructs a confidence region $\mathcal C$
for $P$ under the declared sampling law, an outer effect
region can take the form
\[
 \bigcup_{Q\in\mathcal C}\Theta(Q).
\]
On the event $P\in\mathcal C$, this union contains the true
model's effect, provided the structural class is correct
and the union is evaluated conservatively. That is a
conditional logical implication, not a finite-sample
guarantee for an unspecified estimator. A learned graph
class needs its own retention argument; numerical
approximations to the union need error control as well.

Identification does not by itself control numerical
evaluation. Suppose an identified conditional probability
is the ratio $q=n/d$, with $d>0$. A routine returns
$\widetilde n,\widetilde d$ with genuinely certified
deterministic absolute errors at most $\epsilon$. When
$\widetilde d>\epsilon$ and the probability bounds are
clipped to $[0,1]$, a conservative envelope is
\[
 q\in\left[
 \frac{\max(0,\widetilde n-\epsilon)}{\widetilde d+\epsilon},
 \ \min\left\{1,
 \frac{\widetilde n+\epsilon}{\widetilde d-\epsilon}\right\}
 \right].
\]
The envelope uses monotonicity of a nonnegative ratio.
Correlations between numerator and denominator errors
could tighten it; ignoring those correlations preserves
the outer inclusion. A declared tolerance in an optimizer
is not automatically such a certificate.

Take nominal values
$\widetilde d=10^{-6}$,
$\widetilde n=(3/5)10^{-6}$ and $\epsilon=10^{-7}$.
The envelope is
\[
 q\in[5/11,7/9].
\]
If the comparison risk is exactly $2/5$, the net value
$100(q-2/5)-12$ lies in
$[-72/11,232/9]$ USD. The nominal value is USD 8,
yet the certified numerical envelope crosses zero.
Both nominal masses round to zero at four decimal
places, creating an apparent $0/0$. This is a loss
of numerical information, not proof that the
scientific query is nonidentified.~\citep{exec17_s3_exact}

For comparison, keep the same absolute error but set
$\widetilde d=1/10$ and
$\widetilde n=3/50$. The ratio envelope becomes
\[
 \left[\frac{599999}{1000001},
       \frac{600001}{999999}\right],
\]
with corresponding net values close to USD 8.
The conditioning probability controls how absolute mass
error propagates. In actual data a rare conditioning
event can also cause large sampling uncertainty. The
two effects should be measured separately rather than
combined into an unexplained missing answer.

For the product proposal, this distinction creates
different remedies. A true zero denominator requires
revisiting the target or support assumptions. An
uncertified small positive denominator requires a
numerical diagnostic and a conservative disposition.
Underflow or cancellation may justify rescaling,
log-domain calculation or higher precision followed
by verification. None of these repairs supplies a
missing causal restriction. Conversely, an experiment
that supplies causal information does not repair an
incorrect floating-point implementation.

\section{Choose software by the scientific route it preserves}
\label{exec17:s3-software}

The detailed software comparisons distinguish graphical
adjustment, general intervention identification,
estimation and supplied structural-model simulation.
The synthesis here uses three inspected routes to show
how a product should choose between them; it is not a
claim that these exhaust available identification or
counterfactual software.

DAGitty's inspected adjustment routines reason about a
supplied graph and query. Their useful role is to expose
admissible adjustment choices with the graph's endpoint
meaning intact. They do not learn the graph's scientific
truth from a bank table. The inspected source leaves
some full MAG/PAG validation work as a TODO, and graph
conversion can discard endpoint information. A proposed
integration should therefore validate the admitted
graph class and preserve the original graph alongside
any computational projection. It should reject an
unsupported projection before interpreting an adjustment
result.~\citep{exec13_other_dagitty,exec14_olatent_dagitty}

The inspected y0 high-level API dispatches unconditional
queries to ID and conditioned queries to IDC. Its
default handling turns a particular
\texttt{Unidentifiable} exception into \texttt{None};
strict handling retains that exception. This makes
the exact exception type and query construction
scientifically relevant. A generic missing return
should not erase whether a hedge was obtained, the
input was invalid, an algorithm was unsupported or
an unexpected exception occurred.~\citep{exec13_y0_api,exec13_y0_id,exec13_y0_idc}

The earlier source trace also identifies narrower
concerns. A conditional branch of the inspected query
expression convenience property does not construct
interventions in the same way as the unconditional
branch, although separately stored query fields
remain available. Graph operations can lose an
isolated conditioning node: in the test question
with $A\to Y$ and an isolated $Z$, this may make
computation unavailable even though the stipulated
positive law identifies $P(Y\mid do(A),Z)$.
These are static source findings, not new installed
runtime observations or proofs of mathematical
nonidentification. A suitable integration uses explicit
query fields, checks node preservation and admits a
joint-first conditional fallback only with its
incompleteness clearly stated.~\citep{exec13_y0_api,exec13_y0_graph,exec13_y0_dsl,exec13_y0_utils,exec13_y0_rule2,exec13_y0_separation}

The inspected DoWhy model code distinguishes its
IDIdentifier route from its AutoIdentifier route;
choosing an identification method is not equivalent
to choosing an estimator. Its inspected model call
also assigns a local
\texttt{proceed\_when\_unidentifiable} value without
passing that local at the examined call. The proposal
should test the effective option propagation for
the pinned route and preserve an identification
result separately from later estimation. Repeated
refutation exercises on an estimate cannot establish
assumptions that the underlying observation law
does not identify.~\citep{exec13_dowhy}

For these components, adopting a narrow documented
route behind a scientific interface is preferable
to copying an entire package into a new engine
without a parity argument. A direct port changes
the maintenance and verification obligation:
graph surgery, conditional reduction, expression
semantics and failure witnesses all need matched
cases. The existing source-specific rights and
environment reviews remain necessary inputs.
The present evidence supports a prototype comparison,
not a declaration that one stack has been qualified
for every graph or banking regime.

\section{A bounded prototype can resolve implementation choices}
\label{exec17:s3-design}

The six exact law families supply complementary tests
for a proposed public/generated prototype. The
adjustment family checks population weighting and
the distinction between total and controlled direct
effects. The extended bow family checks conditional
identification without an identified joint. The
coupling family checks refusal to turn identified
experimental marginals into identified personal
attribution. The encouragement family checks LATE,
the observable assignment law and the population
effect bounds. The consistency family checks sharp
extreme completions. The rare-event family checks
whether an otherwise valid formula survives numerical
evaluation. Each family has its own expected scientific
answer; agreement on one does not replace the
others.~\citep{exec17_s3_exact,exec17_s3_decision,exec17_s3_design}

Compare components only on a declared shared domain.
An adjustment-only route should be assessed on
adjustment queries; its refusal of a nonadjustment
query should not count as a false identification.
For supported ID/IDC routes, retain the graph,
intervention and condition sets, positivity domain,
formula or witness, and exact expected result.
For a supplied structural-model simulator, retain
the additional mechanisms and exogenous coupling;
agreement with one selected model must not be scored
as identifying the whole observational class.

A meaningful result distinguishes at least a
point functional, a supported effect set, a valid
nonidentification witness, an unsupported scientific
route, invalid input or support, numerical
unavailability, and an execution failure.
These distinctions guide action. Extra assumptions
can shrink a scientific set; a fuller algorithm may
support an unimplemented route; a software repair
may restore a valid answer; additional observations
may reduce sampling uncertainty. Confusing them
would direct the bank toward the wrong investment.

The initial public prototype can use generated
mechanisms with known responses and permissible
public covariates to make the information
structure realistic. Such covariates do not reveal
unobserved potential responses. Before broader
comparisons, the proposal needs a fixed domain,
version and rights decision for each candidate,
matched output semantics and a separate registered
finite-sample study wherever statistical performance
is claimed. Exact identities and deterministic
error envelopes are a necessary reference, not
that study's outcome.~\citep{exec17_s3_design}

Several limits have a known next step without a
known positive answer. Additional primary-source
review can expand software coverage, but uninspected
candidates cannot yet receive an adoption decision.
Generated comparisons can reveal whether the
adapters preserve the mathematical distinction,
but cannot establish a bank's assignment truth.
Domain work and later intervention evidence can
assess plausible instruments, delivery versions,
interference and harm restrictions. Independent
review, accepted loss thresholds, actual precision
requirements and real costs are later inputs.
The proposal can specify those questions and the
evidence that would answer them while proceeding
with the implementation comparisons whose
requirements are already concrete.
\label{exec17:s3-end}
